% 例 4（平行四边形存在性）：A(0,0), B(3,1), C(1,4)，D 三解：D1(-2,3), D2(2,-3), D3(4,5)
\documentclass[tikz,border=4pt]{standalone}
\usepackage{ctex}
\usepackage{amsmath}
\usetikzlibrary{calc}
\begin{document}
\begin{tikzpicture}[thick, scale=0.55]
  \draw[->, gray] (-4,0) -- (6,0) node[right, font=\footnotesize] {$x$};
  \draw[->, gray] (0,-4) -- (0,6.5) node[above, font=\footnotesize] {$y$};
  \node[font=\footnotesize, below left, gray] at (0,0) {$O$};
  \coordinate (A) at (0,0);
  \coordinate (B) at (3,1);
  \coordinate (C) at (1,4);
  \coordinate (D1) at (-2,3);
  \coordinate (D2) at (2,-3);
  \coordinate (D3) at (4,5);
  \filldraw (A) circle (2pt) node[below left, font=\footnotesize] {$A$};
  \filldraw (B) circle (2pt) node[right, font=\footnotesize] {$B$};
  \filldraw (C) circle (2pt) node[above, font=\footnotesize] {$C$};
  \filldraw[red] (D1) circle (2pt) node[left, red, font=\footnotesize] {$D_1$};
  \filldraw[red] (D2) circle (2pt) node[below, red, font=\footnotesize] {$D_2$};
  \filldraw[red] (D3) circle (2pt) node[above right, red, font=\footnotesize] {$D_3$};
  % 平行四边形 1: A B C D1? 实际情形1: AC, BD 对角线 ⇒ 顺序 A B C D1 围一圈 (ABCD1): 边 AB, BC, CD1, D1A
  \draw[red, thick, dashed] (A) -- (B) -- (C) -- (D1) -- cycle;
  % 情形 2: AB, CD 对角线 ⇒ ACBD2 顺序：A-C-B-D2
  \draw[red, thick, dashed] (A) -- (C) -- (B) -- (D2) -- cycle;
  % 情形 3: BC, AD 对角线 ⇒ ABDC 顺序：A-B-D3-C
  \draw[red, thick, dashed] (A) -- (B) -- (D3) -- (C) -- cycle;
  \node[font=\footnotesize] at (-2,-3) {$D$ 有 $3$ 解};
\end{tikzpicture}
\end{document}
